% SOURCE_AUTHORITY: sga5_fr_workpass.tex, lines 15068--15074 (LF-normalized unit).
% SOURCE_SHA256: 791F4EFFC5E02832D5D77ED03518C8156D6F07E4C8238B03545DB93D883FBB28
Utilicemos ahora los lemas anteriores para reducirnos al caso en que $X$ sea la recta proyectiva.

1\textsuperscript{a} reducción: por inducción noetheriana se puede construir una partición $(U_j)_{1\le j\le r}$ de $X$ por esquemas afines tal que $U_{j+1}$ sea abierto en $X-\bigcup_{i=1}^jU_i$ para $j=1,\ldots,r-1$. Supongamos establecido el teorema para todo par $(U,G)$ formado por un esquema afín $U$ de tipo finito sobre $\mathbb F_p$ y un $\Omega$-haz constructible $G$ sobre $U$; es, por tanto, cierto para $(U_j,F|U_j)$, $j=1,\ldots,r$; gracias al lema 3a), se deduce por inducción descendente sobre $j$ que es cierto para $(\bigcup_{i\ge j}U_i,F|\bigcup_{i\ge j}U_i)$, y por tanto para $(X,F)$, puesto que $X=\bigcup_{i\ge1}U_i$. Basta, pues, demostrar el teorema cuando $X$ es afín.

2\textsuperscript{a} reducción: supongamos que $X$ es afín, de dimensión $n$; existe un morfismo finito $u:X\to\mathbb E^n$ de $X$ al espacio afín de dimensión $n$. Razonemos por inducción sobre $n$. Para $n=0$, $X$ es finito sobre $e$ y el teorema es cierto por el corolario del lema 1. Si $n\ge1$, supongamos establecido el teorema para los esquemas afines de dimensión $n-1$; sean $u:X\to\mathbb E^n$ un morfismo finito y $v:\mathbb E^n\to\mathbb E^1$ un morfismo de proyección, definido a partir de un isomorfismo $\mathbb E^n\simeq\mathbb E^1\times\mathbb E^{n-1}$; al componerlos se obtiene un morfismo $w=v\circ u:X\to\mathbb E^1$ cuyas fibras son esquemas afines de dimensión $n-1$. El lema 1 muestra entonces que la propiedad $(P_{F,w})$ es cierta por la hipótesis de inducción. Gracias al lema 2, se ve que basta establecer el teorema para los pares $(\mathbb E^1,\R^i_!w(F))$. Se reduce así al caso en que $X=\mathbb E^1$ sea la recta afín.

3\textsuperscript{a} reducción: supongamos $X=\mathbb E^1$; sea $i:\mathbb E^1\to\mathbb P^1$ la inmersión de la recta afín en la recta proyectiva $\mathbb P^1$, y sea $Z=\mathbb P^1-i(\mathbb E^1)$ el subesquema cerrado reducido al punto en el infinito de $\mathbb P^1$. Como el teorema es cierto para $(Z,i_!(F)|Z)$, por el lema 3a) equivale establecerlo para $(\mathbb E^1,F)$ o para $(\mathbb P^1,i_!(F))$ (como $i$ es una inmersión abierta, se tiene $F=i_!(F)|\mathbb E^1$).
